\documentclass[a4paper,12pt]{article}
\usepackage[utf8]{inputenc}
\usepackage{geometry}
\usepackage{longtable}
\usepackage{array}
\usepackage{xcolor}
\usepackage{hyperref}
\usepackage{titlesec}
\usepackage{fancyhdr}
\usepackage{enumitem}
\geometry{margin=2cm,headheight=15pt}
\definecolor{brandblue}{RGB}{30,64,175}
\definecolor{wrongred}{RGB}{220,20,60}
\definecolor{correctgreen}{RGB}{34,139,34}
\definecolor{brandgold}{RGB}{181,135,52}
\definecolor{darktext}{RGB}{33,33,33}
\hypersetup{colorlinks=true,linkcolor=brandblue,urlcolor=brandblue}
\pagestyle{fancy}
\fancyhf{}
\fancyhead[L]{\textcolor{brandblue}{\textbf{Approach KAS}}}
\fancyhead[R]{\textcolor{brandblue}{\textbf{UPSC Prelims Test Analysis}}}
\fancyfoot[L]{\textcolor{gray}{\small kas.approachestoias.com}}
\fancyfoot[R]{\textcolor{gray}{\small @ApproachKAS}}
\titleformat{\section}{\Large\bfseries\color{brandblue}}{}{0em}{}[\titlerule]
\titleformat{\subsection}{\large\bfseries\color{darktext}}{}{0em}{}
\begin{document}
\begin{center}
\vspace*{1cm}
{\Huge\bfseries\color{brandblue} UPSC PRELIMS TEST ANALYSIS}\\[0.5cm]
{\Large\color{darktext} Personalized Performance Report}\\[1cm]
\end{center}
\vspace{1cm}
\begin{center}
\fbox{\begin{minipage}{0.9\textwidth}
\centering
\textbf{Student:} Aishwarya Rao \\
\textbf{Test ID:} T004 \\
\textbf{Date:} 2026-10-07 \\
\textbf{Target Exam:} UPSC CSE Prelims 2026
\end{minipage}}
\end{center}
\vspace{1cm}
\section{Executive Summary}
\begin{center}
\begin{tabular}{|c|c|c|c|}
\hline
\textbf{Total Questions} & \textbf{Correct} & \textbf{Wrong} & \textbf{Score} \\
\hline
50 & 28 & 22 & \textcolor{brandblue}{\textbf{56}} \\
\hline
\end{tabular}
\end{center}
\vspace{0.5cm}
\section{Subject-wise Performance}
\begin{longtable}{|l|c|c|c|c|}
\hline
\textbf{Subject} & \textbf{Total} & \textbf{Correct} & \textbf{Wrong} & \textbf{Accuracy} \\
\hline
\endfirsthead
\hline
\textbf{Subject} & \textbf{Total} & \textbf{Correct} & \textbf{Wrong} & \textbf{Accuracy} \\
\hline
\endhead
Polity & 8 & 3 & 5 & \textcolor{wrongred}{37.5\%} \\
\hline
History & 7 & 6 & 1 & \textcolor{correctgreen}{85.7\%} \\
\hline
Geography & 8 & 4 & 4 & \textcolor{brandgold}{50.0\%} \\
\hline
Economy & 7 & 2 & 5 & \textcolor{wrongred}{28.6\%} \\
\hline
Science & 6 & 4 & 2 & \textcolor{brandgold}{66.7\%} \\
\hline
Environment & 6 & 3 & 3 & \textcolor{brandgold}{50.0\%} \\
\hline
Current Affairs & 8 & 6 & 2 & \textcolor{correctgreen}{75.0\%} \\
\hline
\end{longtable}
\vspace{0.5cm}
\section{Mistake Pattern Analysis}
\begin{center}
\begin{tabular}{|l|c|}
\hline
\textbf{Mistake Type} & \textbf{Count} \\
\hline
Partial Knowledge & 6 \\
\hline
Conceptual Gap & 1 \\
\hline
Confusion Similar & 10 \\
\hline
Misinterpretation & 1 \\
\hline
Guesswork & 3 \\
\hline
Factual Error & 1 \\
\hline
\end{tabular}
\end{center}
\section{Focus Areas (Weak Subjects)}
\textbf{Polity (Accuracy: 37.5\%):}\\newline
\textbf{Priority:} \textcolor{wrongred}{HIGH}\\newline
\textbf{Topics to Improve:}\newline
\begin{itemize}
\item Fundamental Rights - Right to Equality (Article 14-18) (1 mistakes)
\item Centre-State Relations - President's Rule (Article 356) (1 mistakes)
\item Fundamental Rights - Right to Freedom of Religion (Article 25-28) (1 mistakes)
\end{itemize}
\textbf{Geography (Accuracy: 50.0\%):}\\newline
\textbf{Priority:} \textcolor{brandgold}{MEDIUM}\\newline
\textbf{Topics to Improve:}\newline
\begin{itemize}
\item Indian Geography – Drainage - River Systems (1 mistakes)
\item Indian Geography – Climate - Climate Regions (1 mistakes)
\item World Geography - Natural Hazards (1 mistakes)
\end{itemize}
\textbf{Economy (Accuracy: 28.6\%):}\\newline
\textbf{Priority:} \textcolor{wrongred}{HIGH}\\newline
\textbf{Topics to Improve:}\newline
\begin{itemize}
\item Banking \& Monetary Policy - RBI Tools (1 mistakes)
\item Budgeting \& Fiscal Policy - Deficit Types (1 mistakes)
\item Government Schemes - DBT and Financial Inclusion (1 mistakes)
\end{itemize}
\textbf{Environment (Accuracy: 50.0\%):}\\newline
\textbf{Priority:} \textcolor{brandgold}{MEDIUM}\\newline
\textbf{Topics to Improve:}\newline
\begin{itemize}
\item Biodiversity - Biodiversity Hotspots (1 mistakes)
\item Environmental Pollution - Air Pollution (1 mistakes)
\item International Agreements - CBD and Nagoya Protocol (1 mistakes)
\end{itemize}
\section{Strength Areas}
\textbf{History}: 85.7\% accuracy - Maintain consistency, focus on advanced questions\\newline
\section{Personalized Recommendations}
\subsection{Recommendation 1: Subject Focus}
\textbf{Priority:} \textcolor{brandgold}{HIGH}\\newline
Focus on Polity - accuracy is 37.5\%. Recommended: Revise NCERT basics and solve 50+ PYQs on weak topics.\\newline
\textbf{Action Items:}\newline
\begin{enumerate}
\item Complete NCERT Polity fundamentals
\item Solve 30 PYQs on Fundamental Rights - Right to Equality (Article 14-18), Centre-State Relations - President's Rule (Article 356), Fundamental Rights - Right to Freedom of Religion (Article 25-28)
\item Take 2 sectional tests on Polity
\end{enumerate}
\subsection{Recommendation 2: Subject Focus}
\textbf{Priority:} \textcolor{brandgold}{MEDIUM}\\newline
Focus on Geography - accuracy is 50.0\%. Recommended: Revise NCERT basics and solve 50+ PYQs on weak topics.\\newline
\textbf{Action Items:}\newline
\begin{enumerate}
\item Complete NCERT Geography fundamentals
\item Solve 30 PYQs on Indian Geography – Drainage - River Systems, Indian Geography – Climate - Climate Regions, World Geography - Natural Hazards
\item Take 2 sectional tests on Geography
\end{enumerate}
\subsection{Recommendation 3: Subject Focus}
\textbf{Priority:} \textcolor{brandgold}{HIGH}\\newline
Focus on Economy - accuracy is 28.6\%. Recommended: Revise NCERT basics and solve 50+ PYQs on weak topics.\\newline
\textbf{Action Items:}\newline
\begin{enumerate}
\item Complete NCERT Economy fundamentals
\item Solve 30 PYQs on Banking \& Monetary Policy - RBI Tools, Budgeting \& Fiscal Policy - Deficit Types, Government Schemes - DBT and Financial Inclusion
\item Take 2 sectional tests on Economy
\end{enumerate}
\subsection{Recommendation 4: Subject Focus}
\textbf{Priority:} \textcolor{brandgold}{MEDIUM}\\newline
Focus on Environment - accuracy is 50.0\%. Recommended: Revise NCERT basics and solve 50+ PYQs on weak topics.\\newline
\textbf{Action Items:}\newline
\begin{enumerate}
\item Complete NCERT Environment fundamentals
\item Solve 30 PYQs on Biodiversity - Biodiversity Hotspots, Environmental Pollution - Air Pollution, International Agreements - CBD and Nagoya Protocol
\item Take 2 sectional tests on Environment
\end{enumerate}
\subsection{Recommendation 5: Mistake Pattern}
\textbf{Priority:} \textcolor{brandgold}{HIGH}\\newline
Your most common mistake type is 'confusion\_similar'. This indicates a specific weakness that needs targeted intervention.\\newline
\textbf{Action Items:}\newline
\begin{enumerate}
\item Create comparison tables for similar concepts
\item Practice 20 questions on distinguishing similar concepts
\item Use mnemonics to remember differences
\end{enumerate}
\subsection{Recommendation 6: Strength Maintenance}
\textbf{Priority:} \textcolor{correctgreen}{LOW}\\newline
Maintain your strong performance in History (85.7\% accuracy).\\newline
\textbf{Action Items:}\newline
\begin{enumerate}
\item Solve advanced PYQs on History
\item Link static topics with current affairs
\item Take 1 sectional test per week on History
\end{enumerate}
\section{Detailed Mistake Analysis}
\subsubsection{Question 1: Polity}
\textbf{Topic:} Fundamental Rights - Right to Equality (Article 14-18)\\newline
\textbf{Correct Answer:} 1 and 2 only\\newline
\textbf{Your Answer:} \textcolor{wrongred}{1, 2 and 3}\\newline
\textbf{Mistake Type:} partial\_knowledge\\newline
\textbf{Explanation:} Statement 1 is correct: Article 14 uses 'persons' not 'citizens', so it applies to all. Statement 2 is correct: Article 15(3) allows special provisions for women and children. Statement 3 is incorrect: Article 16 does not permit reservation based on religion (this was struck down in Indra Sawhney case).\\newline
\vspace{0.5cm}
\subsubsection{Question 3: Polity}
\textbf{Topic:} Centre-State Relations - President's Rule (Article 356)\\newline
\textbf{Correct Answer:} 1 and 3 only\\newline
\textbf{Your Answer:} \textcolor{wrongred}{1 only}\\newline
\textbf{Mistake Type:} conceptual\_gap\\newline
\textbf{Explanation:} Statement 1 is correct: Both Houses must approve within 2 months. Statement 2 is incorrect: Judicial review is NOT barred - the Supreme Court in S.R. Bommai case (1994) established that President's Rule can be reviewed. Statement 3 is correct: Maximum duration is 3 years.\\newline
\vspace{0.5cm}
\subsubsection{Question 5: Polity}
\textbf{Topic:} Fundamental Rights - Right to Freedom of Religion (Article 25-28)\\newline
\textbf{Correct Answer:} Public order, morality, health and other provisions of Part III\\newline
\textbf{Your Answer:} \textcolor{wrongred}{Public order, morality and health only}\\newline
\textbf{Mistake Type:} confusion\_similar\\newline
\textbf{Explanation:} Article 25(1) states that all persons are equally entitled to freedom of conscience, subject to public order, morality and health and to the other provisions of Part III.\\newline
\vspace{0.5cm}
\subsubsection{Question 6: Polity}
\textbf{Topic:} Judiciary - Supreme Court - Original and Appellate Jurisdiction\\newline
\textbf{Correct Answer:} A-1, B-2, C-3, D-4\\newline
\textbf{Your Answer:} \textcolor{wrongred}{A-2, B-1, C-4, D-3}\\newline
\textbf{Mistake Type:} confusion\_similar\\newline
\textbf{Explanation:} Original Jurisdiction - Article 131. Appellate Jurisdiction - Articles 132-134. Advisory Jurisdiction - Article 143. Review Jurisdiction - Article 137.\\newline
\vspace{0.5cm}
\subsubsection{Question 8: Polity}
\textbf{Topic:} Fundamental Duties - List of Fundamental Duties\\newline
\textbf{Correct Answer:} To pay taxes honestly and promptly\\newline
\textbf{Your Answer:} \textcolor{wrongred}{To protect and improve the natural environment}\\newline
\textbf{Mistake Type:} partial\_knowledge\\newline
\textbf{Explanation:} Paying taxes is a legal obligation under various tax laws but is NOT listed as a Fundamental Duty under Article 51A.\\newline
\vspace{0.5cm}
\subsubsection{Question 12: History}
\textbf{Topic:} Modern India – Social Reform - Social Reform Movements\\newline
\textbf{Correct Answer:} 1 and 2 only\\newline
\textbf{Your Answer:} \textcolor{wrongred}{1, 2 and 3}\\newline
\textbf{Mistake Type:} confusion\_similar\\newline
\textbf{Explanation:} Statement 1 is correct: Jyotirao Phule founded Satyashodhak Samaj. Statement 2 is correct: Raja Ram Mohan Roy founded Brahmo Samaj. Statement 3 is incorrect: Swami Vivekananda founded Ramakrishna Mission. Arya Samaj was founded by Swami Dayanand Saraswati.\\newline
\vspace{0.5cm}
\subsubsection{Question 17: Geography}
\textbf{Topic:} Indian Geography – Drainage - River Systems\\newline
\textbf{Correct Answer:} Mahanadi\\newline
\textbf{Your Answer:} \textcolor{wrongred}{Son}\\newline
\textbf{Mistake Type:} confusion\_similar\\newline
\textbf{Explanation:} The Amarkantak Hills is the origin of Narmada, Son, and Johila rivers. The Mahanadi originates from the Siwari hills in Chhattisgarh, not Amarkantak.\\newline
\vspace{0.5cm}
\subsubsection{Question 18: Geography}
\textbf{Topic:} Indian Geography – Climate - Climate Regions\\newline
\textbf{Correct Answer:} 2 and 3 only\\newline
\textbf{Your Answer:} \textcolor{wrongred}{1 and 2 only}\\newline
\textbf{Mistake Type:} misinterpretation\\newline
\textbf{Explanation:} Statement 1 is incorrect: Most peninsular India falls under Tropical Wet and Dry climate. Statement 2 is correct: Western coast falls under Tropical Monsoon (Am). Statement 3 is correct: Himalayan region falls under Highland climate (H).\\newline
\vspace{0.5cm}
\subsubsection{Question 21: Geography}
\textbf{Topic:} World Geography - Natural Hazards\\newline
\textbf{Correct Answer:} Southeast Asia (Indo-Pacific)\\newline
\textbf{Your Answer:} \textcolor{wrongred}{Caribbean Sea}\\newline
\textbf{Mistake Type:} guesswork\\newline
\textbf{Explanation:} The Coral Triangle is in the Western Pacific Ocean around Indonesia, Malaysia, Papua New Guinea, Philippines, Solomon Islands, and Timor-Leste. It contains 76\% of the world's coral species.\\newline
\vspace{0.5cm}
\subsubsection{Question 23: Geography}
\textbf{Topic:} Physical Geography – Oceanography - Ocean Currents\\newline
\textbf{Correct Answer:} Cold current in the Pacific Ocean\\newline
\textbf{Your Answer:} \textcolor{wrongred}{Warm current in the Indian Ocean}\\newline
\textbf{Mistake Type:} confusion\_similar\\newline
\textbf{Explanation:} The Humboldt Current (Peru Current) is a cold, low-salinity ocean current flowing northward along the western coast of South America.\\newline
\vspace{0.5cm}
\subsubsection{Question 24: Economy}
\textbf{Topic:} Banking \& Monetary Policy - RBI Tools\\newline
\textbf{Correct Answer:} Margin Requirements\\newline
\textbf{Your Answer:} \textcolor{wrongred}{Bank Rate}\\newline
\textbf{Mistake Type:} confusion\_similar\\newline
\textbf{Explanation:} Quantitative tools include: Bank Rate, Open Market Operations, CRR, SLR, and Repo/Reverse Repo rates. Margin requirements are qualitative/Selective Credit Control tools.\\newline
\vspace{0.5cm}
\subsubsection{Question 25: Economy}
\textbf{Topic:} Budgeting \& Fiscal Policy - Deficit Types\\newline
\textbf{Correct Answer:} 1, 2 and 3\\newline
\textbf{Your Answer:} \textcolor{wrongred}{1 and 2 only}\\newline
\textbf{Mistake Type:} partial\_knowledge\\newline
\textbf{Explanation:} All three statements are correct. Fiscal Deficit = Total Expenditure - (Revenue Receipts + Non-debt Capital Receipts). Primary Deficit = Fiscal Deficit - Interest Payments. Effective Revenue Deficit = Revenue Deficit - Grants for creation of capital assets.\\newline
\vspace{0.5cm}
\subsubsection{Question 28: Economy}
\textbf{Topic:} Government Schemes - DBT and Financial Inclusion\\newline
\textbf{Correct Answer:} Aadhaar-linked bank accounts and other digital modes\\newline
\textbf{Your Answer:} \textcolor{wrongred}{Cash transfers only}\\newline
\textbf{Mistake Type:} partial\_knowledge\\newline
\textbf{Explanation:} DBT transfers subsidies directly to beneficiaries using Aadhaar-linked bank accounts and other digital payment modes. As of 2024, DBT has saved over ₹3 lakh crore.\\newline
\vspace{0.5cm}
\subsubsection{Question 29: Economy}
\textbf{Topic:} External Sector - Balance of Payments\\newline
\textbf{Correct Answer:} Foreign Direct Investment (FDI)\\newline
\textbf{Your Answer:} \textcolor{wrongred}{Export of software services}\\newline
\textbf{Mistake Type:} confusion\_similar\\newline
\textbf{Explanation:} The Capital Account records: FDI, FPI/FII, external assistance, NRI deposits, and other capital flows. The Current Account records: trade in goods/services, primary income, and secondary income (remittances).\\newline
\vspace{0.5cm}
\subsubsection{Question 30: Economy}
\textbf{Topic:} Poverty \& Unemployment - Measurement Methods\\newline
\textbf{Correct Answer:} There is no official poverty estimation currently\\newline
\textbf{Your Answer:} \textcolor{wrongred}{NITI Aayog}\\newline
\textbf{Mistake Type:} factual\_error\\newline
\textbf{Explanation:} There is NO official poverty estimation currently being done by the Indian government. The last official estimates were based on Tendulkar Committee methodology (2011-12). The government currently uses the Multidimensional Poverty Index (MPI) by NITI Aayog.\\newline
\vspace{0.5cm}
\subsubsection{Question 34: Science}
\textbf{Topic:} Biology – Human Physiology - Immune System\\newline
\textbf{Correct Answer:} 2 and 3 only\\newline
\textbf{Your Answer:} \textcolor{wrongred}{1, 2 and 3}\\newline
\textbf{Mistake Type:} confusion\_similar\\newline
\textbf{Explanation:} Statement 1 is incorrect: Antibodies are produced by B-cells (B-lymphocytes), not T-cells. Statement 2 is correct: Antibodies are proteins called immunoglobulins. Statement 3 is correct: Monoclonal antibodies were used in COVID-19 treatment.\\newline
\vspace{0.5cm}
\subsubsection{Question 36: Science}
\textbf{Topic:} Chemistry - Green Chemistry\\newline
\textbf{Correct Answer:} Paul Anastas and John Warner\\newline
\textbf{Your Answer:} \textcolor{wrongred}{Friedrich Wöhler and Justus von Liebig}\\newline
\textbf{Mistake Type:} guesswork\\newline
\textbf{Explanation:} The 12 Principles of Green Chemistry were developed by Paul Anastas and John Warner in their 1998 book 'Green Chemistry: Theory and Practice'.\\newline
\vspace{0.5cm}
\subsubsection{Question 37: Environment}
\textbf{Topic:} Biodiversity - Biodiversity Hotspots\\newline
\textbf{Correct Answer:} Sundarbans\\newline
\textbf{Your Answer:} \textcolor{wrongred}{Himalayas}\\newline
\textbf{Mistake Type:} confusion\_similar\\newline
\textbf{Explanation:} India has 4 biodiversity hotspots: Western Ghats, Himalayas, Indo-Burma region, and Sundaland (Nicobar Islands). The Sundarbans is a specific ecosystem but NOT classified as a separate biodiversity hotspot.\\newline
\vspace{0.5cm}
\subsubsection{Question 40: Environment}
\textbf{Topic:} Environmental Pollution - Air Pollution\\newline
\textbf{Correct Answer:} 20-30\%\\newline
\textbf{Your Answer:} \textcolor{wrongred}{30-40\%}\\newline
\textbf{Mistake Type:} partial\_knowledge\\newline
\textbf{Explanation:} NCAP launched in 2019 aims to reduce PM2.5 and PM10 concentration by 20-30\% by 2025-26 (later revised to 40\% reduction by 2026).\\newline
\vspace{0.5cm}
\subsubsection{Question 41: Environment}
\textbf{Topic:} International Agreements - CBD and Nagoya Protocol\\newline
\textbf{Correct Answer:} Convention on Biological Diversity (CBD)\\newline
\textbf{Your Answer:} \textcolor{wrongred}{Paris Agreement}\\newline
\textbf{Mistake Type:} confusion\_similar\\newline
\textbf{Explanation:} The Nagoya Protocol is a supplementary agreement to the Convention on Biological Diversity (CBD). It provides a legal framework for fair sharing of benefits from genetic resources.\\newline
\vspace{0.5cm}
\subsubsection{Question 45: Current Affairs}
\textbf{Topic:} Science \& Tech – Recent - Artificial Intelligence\\newline
\textbf{Correct Answer:} AI computing infrastructure\\newline
\textbf{Your Answer:} \textcolor{wrongred}{Cyber defense}\\newline
\textbf{Mistake Type:} guesswork\\newline
\textbf{Explanation:} AIRAWAT is India's AI computing infrastructure, developed by C-DAC under the National AI Mission. It is one of the largest supercomputing platforms for AI workloads in India.\\newline
\vspace{0.5cm}
\subsubsection{Question 48: Current Affairs}
\textbf{Topic:} Government Schemes - Education Reforms\\newline
\textbf{Correct Answer:} 14,500 schools\\newline
\textbf{Your Answer:} \textcolor{wrongred}{5,000 schools}\\newline
\textbf{Mistake Type:} partial\_knowledge\\newline
\textbf{Explanation:} PM SHRI scheme launched in 2022 aims to upgrade 14,500 schools across India to showcase the implementation of NEP 2020.\\newline
\vspace{0.5cm}
\end{document}